\documentclass[12pt]{article}
\usepackage[margin=1in]{geometry}
\usepackage{amsmath}
\usepackage{natbib}
\usepackage{booktabs}
\usepackage{hyperref}

\title{Redemption and Repudiation:\\
Southern State Debts, the Fourteenth Amendment, and the\\
Forging of State Sovereign Immunity, 1868--1890}
\author{George J. Hall and Thomas J. Sargent\thanks{Background memorandum
prepared with Claude for the authors' fiscal-history project. Claims marked
\texttt{[PRIMARY]} in the accompanying \texttt{reconstruction\_debt\_repudiation.bib}
were checked against primary texts (the case reports, the amendment texts, and
Scott's 1893 monograph). Aggregate dollar figures and several page numbers remain
to be pinned; they are collected in Section~\ref{sec:caveats}. This memorandum is
a companion to our notes on the Revolutionary British debts, on the 1840s state
defaults, and on sovereign immunity from \emph{Chisholm} to the 1840s.}}
\date{\today}

\begin{document}

\maketitle

\begin{abstract}
Between the late 1860s and the 1880s the ex-Confederate states ran up large debts
under their Radical Reconstruction governments---mostly by lending the state's
credit to railroads---and then, once conservative Democratic ``Redeemer''
governments took power after the Compromise of 1877, repudiated or drastically
scaled down those debts. The Redeemers justified repudiation on grounds of fraud
and the alleged illegitimacy of the ``carpetbag'' governments that had issued the
bonds; by 1890 the Southern states had voided on the order of \$100--116 million
of principal and escaped a comparable sum in interest. This was a second great
wave of American state repudiation, echoing the defaults of the 1840s, and it must
be distinguished carefully from two adjacent things: the \emph{Confederate} war
debts, which the Fourteenth Amendment's Section~4 had already declared ``illegal
and void''; and the antebellum bank-bond repudiations of the 1840s. Because no
creditor could sue a state---the Eleventh Amendment barred it---bondholders,
including the London houses organized in the Corporation of Foreign Bondholders,
had no judicial remedy, and the British government declined to press their claims.
When bondholders nonetheless tried every litigation workaround, the Supreme Court
closed each one, and in doing so---through \emph{Louisiana v. Jumel}, the
\emph{New Hampshire}/\emph{New York} assignment cases, the Virginia Coupon Cases,
and above all \emph{Hans v. Louisiana} (1890)---forged the modern doctrine of state
sovereign immunity. The Reconstruction repudiations were thus the crucible in which
the constitutional unenforceability of state debt was settled.
\end{abstract}

\section{Introduction: A Second Wave of Repudiation}

Our companion memoranda trace two earlier chapters in the enforceability of
American public debt: how the young republic came to honor the private debts owed
to British merchants (and the treaty-supremacy holding of \emph{Ware v. Hylton});
and how the Eleventh Amendment, provoked by \emph{Chisholm v. Georgia} (1793),
removed the courts from the enforcement of state debts, so that when nine
governments defaulted in the 1840s their creditors could rely only on reputation
\citep{hallsargent_11thamendment, hallsargent_statedefaults1840s}. This memorandum
takes up a second great wave of state repudiation---the Southern repudiations that
followed the end of Reconstruction---because it is there that the several threads
converge: the constitutional unenforceability of state debt, first exposed in the
1840s, was \emph{judicially settled} in the 1880s, in litigation over precisely
these Reconstruction bonds.

\section{The Reconstruction Debt Buildup}
\label{sec:buildup}

Under the biracial, Republican-led governments established during Congressional
Reconstruction, the Southern states borrowed heavily---above all by lending the
state's credit to \emph{railroads}, endorsing or guaranteeing railroad bonds on a
large scale. From 1866 to 1872 the eleven ex-Confederate states are estimated to
have amassed on the order of \$132 million in state debt for railroad subsidies
\citep{white2017republic}. A portion of this borrowing was tainted by genuine
fraud---the Alabama \& Chattanooga endorsement in Alabama, the Bullock-era
endorsements in Georgia, and the Swepson--Littlefield frauds behind North
Carolina's special-tax bonds are the notorious cases \citep{scott1893repudiation}.
That real corruption existed is not in doubt; but, as we note below, the Redeemers'
broader claim that the issuing governments were \emph{illegitimate} was a political
and racial justification, not a neutral finding.

\section{The Fourteenth Amendment, Section 4, and What It Did Not Reach}
\label{sec:amend14}

It is essential to distinguish the Reconstruction debts from the \emph{Confederate}
war debts, because the two are easily conflated and only the latter was addressed
by constitutional amendment. Section~4 of the Fourteenth Amendment (ratified 1868)
provides \citep{amend14sec4}:
\begin{quote}
The validity of the public debt of the United States, authorized by law\,\ldots\
shall not be questioned. But neither the United States nor any State shall assume
or pay any debt or obligation incurred in aid of insurrection or rebellion against
the United States, or any claim for the loss or emancipation of any slave; but all
such debts, obligations and claims shall be held illegal and void.
\end{quote}
The clause did two things. It \emph{guaranteed} the validity of the Union's Civil
War debt---drafted in 1866 by the Joint Committee on Reconstruction out of a fear
that returning Southern representatives might combine to repudiate the war debt or
cut off soldiers' pensions. And it \emph{barred} forever the payment, by the United
States or by any state, of Confederate war debts and of any claim for compensation
for emancipated slaves \citep{amend14sec4}. The Supreme Court's principal reading of
the clause, \emph{Perry v. United States} (1935), later held that its guarantee of
the public debt extends broadly to U.S. obligations, not merely to Civil War bonds
\citep{perry1935}.

The crucial point for our purposes is what Section~4 did \emph{not} reach. It voided
\emph{Confederate} debts---obligations incurred in aid of rebellion between 1861 and
1865. The debts the Redeemers repudiated in the 1870s and 1880s were different
animals: bonds issued \emph{after} the war, in 1868--1874, by the Reconstruction
governments. Their repudiation was neither required nor authorized by Section~4; it
was a separate act, resting on state constitutional amendments and on allegations of
fraud and illegitimacy \citep{scott1893repudiation, ratchford1941american}. A
careful account must keep three repudiations apart: the antebellum bank bonds of the
1840s (treated in our companion memo); the Confederate war debt, voided by
Section~4; and the Reconstruction railroad debt, the subject of this memo.

\section{Redemption and Repudiation, State by State}
\label{sec:states}

The mechanism was almost everywhere the same: a Redeemer constitutional convention
or amendment declared specified Reconstruction bonds fraudulent, unconstitutional,
or void, or a ``funding''/``settlement'' act exchanged them for new bonds at a
fraction of face value. Table~\ref{tab:states} summarizes the debts outstanding when
repudiation sentiment crystallized (from Scott's 1893 tabulation) and the outcome.

\begin{table}[htbp]
\centering
\caption{Southern state debts and their repudiation/readjustment}
\label{tab:states}
\small
\begin{tabular}{l r p{7.6cm}}
\toprule
State & Debt at repudiation & Mechanism and outcome \\
\midrule
North Carolina & \$15{,}000{,}000 & Most aggressive: \$16.24M of ``special-tax'' railroad bonds (Swepson--Littlefield fraud); acts repealed, 1880 amendments barred payment \\
Louisiana & \$22{,}500{,}000 & 1874 Funding Act: ``consolidated bonds'' at 60c on the dollar; \$8.09M voided as beyond the constitutional limit; 1879 Debt Ordinance cut the coupon. Source of the leading Eleventh-Amendment cases \\
Virginia & \$45{,}000{,}000 & Largely \emph{antebellum} debt; Riddleberger Act (1882) reset it at $\sim$\$21M in 3\% 50-year bonds, repudiating accrued interest and the West Virginia ``third'' \\
Tennessee & \$43{,}950{,}000 & 1883 Settlement Act scaled the debt to $\sim$50c on the dollar at 3\% (readjustment) \\
Alabama & \$11{,}345{,}000 & 1875 constitution debt limits; 1876 settlement act scaled down the railroad-aid debt \\
Georgia & \$11{,}135{,}500 & Bullock-era endorsements (Brunswick \& Albany RR) declared null and void; 1877 constitution forbade paying them \\
South Carolina & \$15{,}850{,}000 & 1873 Consolidation Act at 50c on the dollar; \$5.9M of ``conversion bonds'' declared void \\
Arkansas & \$8{,}600{,}000 & Railroad-aid, levee, and Holford bonds held void; Fishback Amendment (1884) barred all payment \\
Florida & \$1{,}850{,}000 & Repudiated $\sim$\$8M of bank and railroad-aid bonds (with interest) \\
\bottomrule
\end{tabular}

\smallskip
\footnotesize Source: debt totals from \citet[215]{scott1893repudiation} (Florida
detail at p.~54). Virginia and Tennessee totals largely reflect \emph{antebellum}
internal-improvement debt carried and readjusted after the war. Mississippi is
omitted: its famous repudiation (the Union and Planters Bank bonds) was
\emph{antebellum} (1842), reaffirmed in the 1875 and 1890 constitutions, not a
Reconstruction-debt repudiation. Texas ran up little debt and did not
significantly repudiate.
\end{table}

\subsection{Virginia: the Readjusters as a counterpoint}

Virginia deserves emphasis because it complicates the simple ``Redeemers
repudiated'' story. Virginia's debt was overwhelmingly \emph{antebellum}
internal-improvement debt, swollen to more than \$45 million by the 1871 Funding
Act, which made bond coupons receivable for state taxes---so that within two years
nearly half the state's revenue arrived as coupons, starving the new public
schools \citep{ev_debtcontroversy}. The scaling-down was carried not by conservative
``Redeemers'' but \emph{against} them: the biracial \emph{Readjuster} coalition led
by the former Confederate general William Mahone, which won the legislature in 1879
and the governorship in 1881, and whose Riddleberger Act (1882) reset the debt at
about \$21 million in 3\% bonds to fund schools and services. The conservative
``Funders'' had insisted on paying in full to preserve the state's honor and credit
\citep{moore1974twopaths, maddex1970virginia, ev_readjuster}. Virginia thus shows
that debt scaling could be a \emph{reform} program, not only a Lost-Cause one.

\subsection{A note on the repudiators' rhetoric}

The Redeemer governments justified repudiation as the cleansing of fraud committed
by illegitimate ``black Republican,'' ``scalawag,'' and ``carpetbagger''
governments. Some of the underlying fraud was real. But the ``illegitimacy''
argument was itself part of the Lost-Cause project: the Reconstruction governments
were the lawful, biracial governments of their states, and the racial framing of
repudiation as the repudiation of \emph{their} debts was a political act. A book
treatment should present ``illegitimate carpetbag debt'' as the \emph{repudiators'
characterization}, not as an established fact \citep{white2017republic,
woodward1951origins}.

\section{Foreign Bondholders Without a Remedy}
\label{sec:foreign}

As in the 1840s, much of the repudiated paper was held abroad, chiefly in London.
The \emph{Corporation of Foreign Bondholders}, organized in London in 1868, took up
the creditors' cause and published lists of the defaulted Southern debts. But the
creditors had essentially no recourse. The British Foreign Office declined to press
their claims, reasoning---correctly---that even the \emph{federal} government of the
United States had no power to compel a state to pay \citep{mcgrane1935foreign}. The
barrier was the Eleventh Amendment and the sovereign immunity of the states: a
foreign or out-of-state creditor simply could not bring a defaulting state to court.
Enforcement, as in the 1840s, could come only through reputation and the market's
memory---and here even that lever was weak, because the repudiating states largely
did without new foreign borrowing.

\section{The Litigation That Forged Sovereign Immunity}
\label{sec:litigation}

The Reconstruction repudiations produced the litigation in which the modern law of
the Eleventh Amendment was made. Congress had granted the federal courts general
federal-question jurisdiction in 1875, just as the Redeemers were voiding bonds;
bondholders poured into federal court armed with the Contract Clause; and the
Supreme Court, unwilling either to abandon contract doctrine or to provoke wholesale
state defiance of its judgments, resolved the tension by \emph{expanding state
sovereign immunity} \citep{orth1987judicial}. It closed the workarounds one by one:
\begin{itemize}
  \item \emph{Louisiana v. Jumel} (1883): a suit to compel Louisiana's officers to
    pay the repudiated consolidated-bond interest was barred---the money in the
    treasury is the state's, and the state is ``the real and only party'' in
    interest \citep{jumel1883}.
  \item \emph{New Hampshire v. Louisiana} and \emph{New York v. Louisiana} (1883):
    bondholders assigned their Louisiana bonds to their home states so those states
    could sue in the Supreme Court's original jurisdiction (state-versus-state suits
    being unbarred); the Court saw through the device and dismissed
    \citep{nhlouisiana1883}.
  \item The Virginia Coupon Cases: after Virginia repudiated its tax-receivable
    coupons, the Court first upheld the state's substitute remedy (\emph{Antoni v.
    Greenhow}, 1883), then allowed a suit against an officer who seized property
    under a void law as a private wrongdoer (\emph{Poindexter v. Greenhow}, 1885),
    then cabined that opening by holding that a suit to make officers honor the
    state's \emph{contract} was really a suit against the state (\emph{Ex parte
    Ayers}, 1887) \citep{antoni1883, poindexter1885, ayers1887}.
  \item \emph{Hans v. Louisiana} (1890): the capstone. A Louisiana citizen sued his
    own state for the repudiated consolidated-bond interest, arguing the Contract
    Clause; the Court held that a state may not be sued in federal court even by its
    own citizens, extending immunity beyond the Amendment's literal text
    \citep{hans1890}.
\end{itemize}
As John Orth argues, \emph{Hans} and its companions are best explained by the
post-Reconstruction political settlement: the Court used the Eleventh Amendment to
shield the Southern states from their creditors and ease sectional reconciliation
\citep{orth1987judicial, purcell2003hans, gibbons1983eleventh}. In James Ely's
phrase, \emph{Hans} ``can best be understood as part of the Supreme Court's refusal,
on claimed jurisdictional grounds, to confront the widespread repudiation of bonds
by southern states.'' The constitutional unenforceability of state debt that the
1840s had \emph{revealed} was, in the 1880s, \emph{entrenched}.

\section{Connections to the Book and the Companion Memoranda}
\label{sec:connections}

Several links to the other memoranda seem worth carrying into the manuscript.
\begin{enumerate}
  \item \textbf{The arc of state sovereign immunity closes here.} Our
    Eleventh-Amendment memo runs from \emph{Chisholm} (1793) to the 1840s;
    \emph{Hans v. Louisiana}, which appears there as the doctrine's capstone,
    \emph{arose from this episode}---Louisiana's repudiation of its 1874
    consolidated bonds. This memo supplies the factual matrix for that case and its
    companions. The two memos are naturally read as one arc.
  \item \textbf{Two waves of state repudiation.} The 1840s and the 1870s--80s are
    the two great waves, and they rhyme: heavy borrowing (canals/banks, then
    railroads), a shock (the Panic of 1837, then the collapse of Reconstruction and
    the Panic of 1873), default and repudiation, foreign creditors without recourse,
    and a constitutional response (the debt-limit constitutions of 1842--52, then the
    sovereign-immunity cases). The reputational logic of our 1840s memo applies
    again---with the twist that the repudiators here mostly did not seek to
    re-borrow abroad, so market discipline bit less.
  \item \textbf{Federal versus state credit.} As with the Revolutionary British
    debts, a \emph{constitutional} obligation could be enforced (Section~4 protects
    the \emph{Union} debt, vindicated in \emph{Perry}), while a state's own debt
    could not (the Eleventh Amendment). The contrast between enforceable federal
    credit and unenforceable state credit, a theme of all four memos, is at its
    sharpest here: the same Fourteenth Amendment that guaranteed the federal debt sat
    beside an Eleventh Amendment that let the states repudiate.
  \item \textbf{Virginia recurs.} Virginia appears in three of the memos---its
    Revolutionary British debts, and its post-war Readjuster settlement---a
    convenient through-line if the book wants one state's fiscal biography.
\end{enumerate}

\section{Caveats and Verification Notes}
\label{sec:caveats}

\begin{itemize}
  \item \textbf{Aggregate repudiation figures.} The ``$\sim$\$116 million of
    principal repudiated and $\sim$\$150 million of interest escaped by 1890'' figure
    is attributed to \citet{white2017republic} but not pinned to a page; an
    alternative ``\$75M principal / \$300M interest'' circulates as ``from Scott''
    but we could not locate it in Scott's text. Confirm against \citet{mcgrane1935foreign}
    (pp.~282--381) and \citet{ratchford1941american} before printing a total.
  \item \textbf{The \$132 million railroad-subsidy figure} is attributed to
    \citet{white2017republic}; verify the page.
  \item \textbf{Per-state amounts} in Table~\ref{tab:states} are debts
    \emph{outstanding} at repudiation (Scott, p.~215), not net amounts repudiated;
    Alabama's scaled-down figure in particular was not firmly verified. Scott page
    numbers other than p.~54 and p.~215 are approximate (from OCR) and should be
    checked.
  \item \textbf{Virginia's exact Riddleberger principal} (often given as
    \$21,035,377.15) is confirmed only as ``a little more than \$21 million'';
    confirm the exact figure in \citet{moore1974twopaths} or the statute.
  \item \textbf{Section 4 drafting motive.} The claim (sometimes made) that concern
    over \emph{British-held} Confederate bonds motivated Section~4 could not be
    confirmed and is not asserted here.
  \item \textbf{Orth page numbers} for the repudiation-cases thesis are given as a
    chapter locator ($\sim$pp.~58--89), not a pinpoint; the Ely quotation should be
    verified against its source before quoting.
\end{itemize}

\bibliographystyle{aer}
\bibliography{reconstruction_debt_repudiation}

\end{document}
